\section{The comparison measure and its elementary bounds}
\label{appa-measure}

This is the analogue of source Appendix~A, pp.~21--26.  The measure and
the numerical enclosures are \textbf{[Formal analogue: reused source
input]}; their defining algebra is displayed below.  No new search for
an equilibrium measure or independent verification of the source's
684 partition intervals is claimed.

\subsection{The sixteen intervals}

Let $d\omega_{[a,b]}(t)=dt/(\pi\sqrt{(t-a)(b-t)})$ on $(a,b)$ and set
$\rho=\sum_{j=1}^{16}c_j\omega_{[a_j,b_j]}$.  The following integers,
each divided by $10^{12}$, reproduce source Table~1:
\begin{center}
\begin{tabular}{r|r|r|r}
$j$&$10^{12}a_j$&$10^{12}b_j$&$10^{12}c_j$\\\hline
1&3906748086&8992695531&10515596180\\
2&2312248264&15340997855&29471737793\\
3&1402286665&25730180724&42934365099\\
4&881725356&41909578246&58204231966\\
5&578197906&65851089563&69037621310\\
6&396324613&99481037884&78873099189\\
7&283911191&144325727458&84856120711\\
8&212206188&201105762729&88396082127\\
9&165097686&269345996903&88303382125\\
10&133347132&347089554156&85472321255\\
11&111522114&430806704415&78899184238\\
12&96349355&515561896511&70353471918\\
13&85815639&595448778546&58838976615\\
14&78667711&664241383483&44421321106\\
15&74129565&716160577112&30462865791\\
16&71741310&746637295669&5959622577
\end{tabular}
\end{center}
Direct rational arithmetic gives $\sum c_j=37/40$,
$0<a_{16}<\cdots<a_1<b_1<\cdots<b_{16}<2$, and
$b_j-a_j>1/225$.  With $m=(a+b)/2$, the substitution
$u=m+(b-a)\cos\theta/2$ gives
\begin{equation}
 U_{\omega_{[a,b]}}(t)=
 \begin{cases}
 \log((b-a)/4),&a\le t\le b,\\
 \displaystyle\log\frac{|t-m|+\sqrt{(t-a)(t-b)}}2,&t\notin[a,b].
 \end{cases}
 \label{appa-arcsine}
\end{equation}
For the exterior branch, differentiate the defining integral and
normalize by $U(t)-\log|t|\to0$; on the interval, the sine-product
identity for $\cos\phi-\cos\theta$ gives the constant.  This proves
the algebra of (A.1).

Put $S_j=\sum_{i\le j}c_i$, $S_0=0$.  An outer interval's potential
is constant on every inner interval, so the self and cross terms combine
as $c_j^2+2c_jS_{j-1}=S_j^2-S_{j-1}^2$.  Hence (A.2) becomes
\begin{equation}
 I(\rho)=\sum_{j=1}^{16}(S_j^2-S_{j-1}^2)
                          \log\frac{b_j-a_j}{4}.
 \label{appa-energy}
\end{equation}

\subsection{Elementary enclosures}

For $1\le x\le2$, $z=(x-1)/(x+1)$, the logarithmic series gives (A.3):
\[
 0\le\log x-2\sum_{k=0}^{m-1}\frac{z^{2k+1}}{2k+1}
       \le\frac{2z^{2m+1}}{(2m+1)(1-z^2)}.
\]
Reduce a general argument as $x=2^er$, $1\le r<2$.  For
$0\le z\le1/2$, the alternating series gives (A.4):
\[
 \left|\arctan z-\sum_{k=0}^{m-1}
                  \frac{(-1)^kz^{2k+1}}{2k+1}\right|
       \le\frac{z^{2m+1}}{2m+1}.
\]
Use $\arctan x=\pi/2-\arctan(1/x)$ for $x>1$ and
$\arctan x=2\arctan(x/(1+\sqrt{1+x^2}))$ for $0\le x\le1$;
enclose $\pi$ through $16\arctan(1/5)-4\arctan(1/239)$.
The source chooses $m=64$ for logarithms, $m=80$ for arctangents,
and adjacent dyadic multiples of $2^{-144}$ for square roots.
The reduced series errors are below $2^{-144}$; interval arithmetic
must also propagate the rounding errors through each expression.

\subsection{The potential inequality on the half-line}

Integrating the field in \eqref{s6-field} gives (A.5):
\begin{align}
 V(t)&=\log(1+t)-6\alpha\log(t+\alpha^2)-2+12\alpha\notag\\
 &\quad+2\sqrt t\left(\pi+\arctan\frac1{\sqrt t}
                                -6\arctan\frac\alpha{\sqrt t}\right),
 \label{appa-field-explicit}
\end{align}
with $V(0)=-12\alpha\log\alpha-2+12\alpha$.
For $\Phi(y)=V(y^2)$,
\[
 \Phi'(y)=2\left(\pi+\arctan(1/y)-6\arctan(\alpha/y)\right),
 \qquad \Phi''(y)=-\frac2{1+y^2}+\frac{12\alpha}{\alpha^2+y^2}.
\]
The latter changes sign at $y^2=711/880$, while
$\Phi'(0+)=-3\pi$ and $\Phi'(\infty)=2\pi$.  Thus $V$ has one minimum.
The source encloses it between
$q_-=59205077/10^{10}$ and $q_+=59205079/10^{10}$ by opposite
signs of $\Phi'(\sqrt{q_\pm})$.
On this small interval a lower bound, as in (A.6), is
\begin{align*}
 V_*={}&\log(1+q_-)-6\alpha\log(q_++\alpha^2)-2+12\alpha\\
 &+2\sqrt{q_+}\left(\pi+\arctan\frac1{\sqrt{q_+}}
                              -6\arctan\frac\alpha{\sqrt{q_-}}\right).
\end{align*}
The parenthesis is negative, which determines the direction of the
square-root bound.  Nesting makes $U_\rho$ nonincreasing up to $a_1$,
constant on $[a_1,b_1]$, and nondecreasing thereafter.  Thus (A.7) is
\[
 B(l,r)=2\max\{U_\rho(l),U_\rho(r)\}
 -\begin{cases}V(r),&r\le q_-,\\V(l),&l\ge q_+,\\V_*,&\text{otherwise},\end{cases}
 \qquad 2U_\rho(t)-V(t)\le B(l,r)\quad(l\le t\le r).
\]
For completeness, the partition rule (A.8) is
\[
 A_0=0,\quad A_j=a_{17-j}\ (1\le j\le16),\quad
 A_{17}=q_-,\quad A_{18}=q_+,\quad
 A_{18+j}=b_j\ (1\le j\le16),\quad A_{35}=2,
\]
where a source Table~2 row $(j,d,k)$ specifies
\[
 \left[A_j+(A_{j+1}-A_j)\frac{k}{2^d},\quad
       A_j+(A_{j+1}-A_j)\frac{k+1}{2^d}\right].
\]
Use exactly the row list in source Table~2, pp.~24--25.  Its claimed
rational certificate (A.9) is
\[
 B(l,r)<-\frac{6645002}{10^6}<-\frac{1329}{200}
\]
on every listed interval.  The tail estimate \eqref{s6-tail} treats
$t\ge2$.  Because both $\rho$ and $V$ are unchanged, these are the same
numerical inputs in the $\zeta(9)$ draft, not new checks.

\subsection{Energy and the norm constant}

Applying the stated logarithmic enclosures to \eqref{appa-energy} and
$C_*$ gives the source input (A.10):
\[
 -\frac{2126593445148}{10^{12}}<I(\rho)
              <-\frac{2126593445147}{10^{12}},\qquad
 \frac{2653035990340}{10^{12}}<C_*
              <\frac{2653035990341}{10^{12}}.
\]
Since $\lambda M_0=-49173/8000$, the directions of the bounds give
\[
 \lambda M_0-I(\rho)+C_*<-
 \frac{1366995564511}{10^{12}}<-
 \frac{2733991}{2000000}=U.
\]
Finally, (A.11) encloses the source's norm exponent
\[
 \kappa=\frac{2\pi}{\sqrt3}+\log\frac43+\frac9{20}\log18
 +\frac{37}{40}\left(1+\log\frac{1600}{111}\right)
 +\frac{37}{20}\log\frac{43+4\sqrt{114}}5
\]
by
\[
 \frac{2773171335773}{200000000000}<\kappa
 <\frac{6932928339433}{500000000000}<\frac{6933}{500}.
\]
These rational enclosures remain source input.  The reason the same
leading norm exponent is available for the fixed-shape analogue is that
the lower bound from the $\ell=1$ term of $w_9$ changes powers and
constants, whereas its exponential factor remains $e^{-2\pi y}$;
the separate norm manipulation belongs to Appendix~C.
